\documentclass[preprint,12pt]{elsarticle}

%-------------------------------------------------
% Mathematics
%-------------------------------------------------
\usepackage{amsmath}
\usepackage{amssymb}
\usepackage{amsfonts}
\usepackage{bm}
\usepackage{mathtools}
\usepackage[a4paper,margin=1in]{geometry}


\DeclareMathOperator{\tr}{tr}
\DeclareMathOperator{\sym}{sym}

%-------------------------------------------------
% Figures and Tables
%-------------------------------------------------
\usepackage{graphicx}
\usepackage{booktabs}
\usepackage{multirow}
\usepackage{array}
\usepackage{subcaption}

%-------------------------------------------------
% Units
%-------------------------------------------------
\usepackage{siunitx}

%-------------------------------------------------
% Colors
%-------------------------------------------------
\usepackage{xcolor}

%-------------------------------------------------
% Hyperlinks (Load Near the End)
%-------------------------------------------------
\usepackage[hidelinks]{hyperref}

\journal{Journal Name}

\begin{document}

\begin{frontmatter}

%%%%%%%%%%%%%%%%%%%%%%%%%%%%%%%%%%%%%%%%%%%%%%%%%%
% Title
%%%%%%%%%%%%%%%%%%%%%%%%%%%%%%%%%%%%%%%%%%%%%%%%%%

\title{A Total-Lagrangian \texorpdfstring{$C^1$}{C1} Bogner--Fox--Schmit Formulation for Finite-Strain Mindlin Form-II Strain-Gradient Elasticity}

%%%%%%%%%%%%%%%%%%%%%%%%%%%%%%%%%%%%%%%%%%%%%%%%%%
% Authors
%%%%%%%%%%%%%%%%%%%%%%%%%%%%%%%%%%%%%%%%%%%%%%%%%%

\author{}
\address{}

%%%%%%%%%%%%%%%%%%%%%%%%%%%%%%%%%%%%%%%%%%%%%%%%%%
% Abstract
%%%%%%%%%%%%%%%%%%%%%%%%%%%%%%%%%%%%%%%%%%%%%%%%%%

\begin{abstract}

A displacement-based total-Lagrangian finite element formulation is developed for geometrically nonlinear Mindlin Form-II strain-gradient elasticity. The formulation employs the Green--Lagrange strain and its material gradient as constitutive variables together with a Saint--Venant--Kirchhoff hyperelastic model and an isotropic sixth-order Mindlin strain-gradient constitutive law. A bicubic $C^1$ Bogner--Fox--Schmit finite element satisfies the continuity requirements of the displacement-only weak form, and an exact consistent tangent matrix yields quadratic convergence in the Newton--Raphson solver. The formulation is verified against the Shekarchizadeh et al. (2022) two-dimensional simple-shear benchmark, recovering the analytical displacement profiles with observed convergence rates close to four in $L^2$, three in $H^1$, and two in the strain-gradient measure. Cantilever bending is then analyzed under finite deformation to assess mesh convergence, internal length-scale effects, higher-order boundary conditions, and the sensitivity to the five isotropic Mindlin constants. The results show that strain-gradient contributions increase apparent cantilever stiffness by roughly 5\%, reduce tip displacement by approximately 4.3--4.7\%, and regularize localized strain concentrations near the clamped boundary. The formulation provides a displacement-only finite element framework for Mindlin Form-II strain-gradient elasticity undergoing finite deformation, operating entirely on a single primal displacement field.

\end{abstract}

%%%%%%%%%%%%%%%%%%%%%%%%%%%%%%%%%%%%%%%%%%%%%%%%%%
% Keywords
%%%%%%%%%%%%%%%%%%%%%%%%%%%%%%%%%%%%%%%%%%%%%%%%%%

\begin{keyword}
Finite-strain elasticity \sep Strain-gradient elasticity \sep Mindlin Form-II elasticity \sep Total-Lagrangian formulation \sep Bogner--Fox--Schmit element \sep $C^1$ finite elements \sep Consistent tangent matrix \sep Newton--Raphson method
\end{keyword}

\end{frontmatter}


\section{Introduction}

\label{sec:introduction}

Classical elasticity describes a body through local stress--strain
relations. It remains valid when structural dimensions are significantly
larger than the characteristic scale of the material microstructure.
However, classical models become inadequate when deformation gradients vary
rapidly near supports, interfaces, cracks, or geometric discontinuities, or
when the mechanical response exhibits intrinsic size dependence.
Strain-gradient elasticity addresses these limitations by introducing
higher-order deformation gradients and internal material length scales
\cite{Aifantis1999,AskesAifantis2002,LamYangChongWangTong2003,AskesAifantis2011}.

The foundations of higher-order elasticity were established by Toupin
\cite{Toupin1962} and Mindlin and Tiersten \cite{MindlinTiersten1962} for
materials with couple stresses, followed by Mindlin's formulations of
microstructural effects \cite{Mindlin1964} and second strain gradients with
surface energy \cite{Mindlin1965}. In first strain-gradient theory,
Mindlin and Eshel \cite{MindlinEshel1968} formulated strain-energy
densities dependent on the infinitesimal strain tensor and its gradient.
Germain \cite{Germain1973} subsequently formulated the principle of
virtual power for strain-gradient media, identifying the dual higher-order
internal stress measures and conjugate boundary tractions upon which
displacement-based variational formulations rely.

The physical interpretation of internal material lengths has been
addressed across multiple scales and mechanisms. Aifantis \cite{Aifantis1999}
related gradient terms to size-dependent elastic and plastic responses, while
Fleck and Hutchinson \cite{FleckHutchinson1993,FleckMullerAshbyHutchinson1994,FleckHutchinson1997}
incorporated strain gradients into phenomenological descriptions of
nonuniform plastic deformation. Lam, Yang, Chong, Wang, and Tong
\cite{LamYangChongWangTong2003} combined experiments and theory to examine
elastic size effects in small-scale beams, and Yang, Chong, Lam, and Tong
\cite{YangChongLamTong2002} introduced a couple-stress strain-gradient
formulation governed by a single internal length scale.

Recent investigations have extended strain-gradient concepts to
architected metamaterials, coupled multiphysics, and generalized continua.
Desai and Sidhardh \cite{DesaiSidhardh2025} reviewed contemporary
strain-gradient theories, computational schemes, and length-scale
applications. Rodopoulos and Karathanasopoulos
\cite{RodopoulosKarathanasopoulos2024} evaluated thermally and mechanically
induced strain-gradient fields in architected materials and beam
structures. Gupta, Barak, and Das \cite{GuptaBarakDas2024} analyzed a
size-dependent nonlocal thermo-piezo-photoelectric medium. In wave studies,
Das et al. \cite{DasDuttaCraciunSurBarakGupta2024} investigated
size-dependent surface-wave interaction in a micropolar thermoelastic medium,
Das et al. \cite{DasDuttaGuptaSinghalBarakAlmohsen2025} examined wave
reflection in prestressed microstructured solids with dual porosity, and
Das et al. \cite{DasSurGuptaDuttaSinghalKumar2025} analyzed coupled
responses in a size-dependent porous piezoelectric medium. Although these
multiphysics formulations differ from the purely mechanical setting
investigated here, they demonstrate the widespread adoption of
length-dependent generalized continua for microstructured media.

The finite-element treatment of strain-gradient elasticity is governed by
the regularity required by the variational weak form. Because a
displacement-based strain-gradient weak statement contains second
derivatives of displacement, it requires $H^2$-conforming
($C^1$-continuous) approximations. Bogner, Fox, and Schmit
\cite{BognerFoxSchmit1966} developed a bicubic Hermite rectangular element
providing global $C^1$ inter-element continuity on regular quadrilateral
meshes. Shu, King, and Fleck \cite{ShuKingFleck1999} discussed the
trade-offs between $C^1$ displacement interpolation and mixed
multi-field formulations. Amanatidou and Aravas \cite{AmanatidouAravas2002}
introduced mixed finite-element formulations to circumvent the requirement
of globally $C^1$-continuous displacement fields. Subsequent developments
in displacement-based and mixed formulations were conducted by Zervos,
Papanastasiou, and Vardoulakis \cite{ZervosPapanastasiouVardoulakis2001},
Askes and Aifantis \cite{AskesAifantis2006}, Askes and Guti\'errez
\cite{AskesGutierrez2006}, Askes, Morata, and Aifantis
\cite{AskesMorataAifantis2008}, and Zervos, Papanicolopulos, and
Vardoulakis \cite{ZervosPapanicolopulosVardoulakis2009}.

Recent investigations have addressed these discretization challenges
through diverse numerical strategies. Choi, Lee, and Sim
\cite{ChoiLeeSim2023} developed mixed elements based on superconvergent
patch recovery, while Papanicolopulos \cite{Papanicolopulos2024} proposed
mixed elements with shared degrees of freedom. Chirkov, Nazarenko, and
Altenbach \cite{ChirkovNazarenkoAltenbach2024} applied a mixed formulation to
plane crack problems in strain-gradient elasticity, and subsequently examined
edge-cracked beam bending \cite{ChirkovNazarenkoAltenbach2025}. Bacciocchi and
Fantuzzi \cite{BacciocchiFantuzzi2023} formulated higher-order Hermite
mappings for gradient-elastic plates on non-rectangular domains, and later
extended the formulation to distorted domains \cite{BacciocchiFantuzzi2025}.
Sarar, Yildizdag, Fabbrocino, and Abali
\cite{SararYildizdagFabbrocinoAbali2025} conducted comparative studies of
Hermite, mixed, and isogeometric discretizations for strain-gradient problems.
Across these studies, the continuity of the chosen shape functions and the
enforcement of higher-order boundary constraints represent the primary
determinants of numerical accuracy and stability.

A compatible $C^1$ displacement-based formulation remains advantageous when
the strain-gradient field must be evaluated directly from a single primal
displacement field without introducing auxiliary degrees of freedom.
Beheshti \cite{Beheshti2017} demonstrated the efficacy of Hermite-type
elements for plane-strain strain-gradient elasticity. Andreaus, dell'Isola,
Giorgio, Placidi, Lekszycki, and Rizzi
\cite{AndreausDellIsolaGiorgioPlacidiLekszyckiRizzi2016} showed that
second-gradient contributions can significantly alter localized stress and
strain fields even when the macroscopic response appears unaffected. The
foundational nonlinear finite-element frameworks of Hughes \cite{Hughes1987},
Belytschko, Liu, and Moran \cite{BelytschkoLiuMoran2000}, Bathe
\cite{Bathe2006}, and Simo and Hughes \cite{SimoHughes1998} provide the
total-Lagrangian kinematics and consistent linearization required for robust
nonlinear implementation.

Finite deformation introduces further requirements: the strain measure must be
objective, the stresses must be work-conjugate to the chosen kinematic
variables, and the tangent stiffness matrix must be the exact derivative of
the nonlinear residual. Standard treatments of finite elasticity are given by
Holzapfel \cite{Holzapfel2000}, Bonet and Wood \cite{BonetWood2008}, and
Ogden \cite{Ogden1997}. Javili, dell'Isola, and Steinmann
\cite{JaviliDellIsolaSteinmann2013} treated geometrically nonlinear
higher-gradient elasticity with energetic boundaries. Abali, M\"uller,
and Eremeyev \cite{AbaliMullerEremeyev2015} presented a computational
strain-gradient framework with geometric nonlinearities, while Abali,
M\"uller, and dell'Isola \cite{AbaliMullerDellIsola2017} developed
higher-gradient formulations based on action principles. In multiscale modeling,
Rodrigues Lopes and Andrade Pires \cite{RodriguesLopesAndradePires2022} developed
a finite-strain second-order homogenization scheme, and Schmidt, Kr\"uger,
Keip, and Hesch \cite{SchmidtKrugerKeipHesch2022} formulated computational
homogenization for higher-order continua. Daneshyar, Sotoudeh, and
Ghaemian \cite{DaneshyarSotoudehGhaemian2022} developed a scaled-boundary
finite element method for higher-order continua with dispersive wave
propagation. In addition, Shakeri, Ghaffari, Thompson, and Brown
\cite{ShakeriGhaffariThompsonBrown2024} identified numerical stability
issues that arise when finite-strain quantities are evaluated near
the reference configuration.

Rigorous verification becomes essential when finite kinematics and
strain-gradient formulations are coupled. Shekarchizadeh, Abali, and Bersani
\cite{ShekarchizadehAbaliBersani2022} developed an analytical simple-shear
benchmark incorporating periodic constraints, higher-order boundary conditions,
and both displacement- and traction-controlled loading cases. Because this
benchmark provides exact analytical solutions for the full displacement field
and its spatial derivatives, it is well suited to verifying the small-strain
limiting behavior of a $C^1$ strain-gradient formulation.

Despite these advances, there remains a need for a displacement-based,
total-Lagrangian formulation that preserves the full kinematic description of
higher-order elasticity while accommodating finite strains, general
higher-order boundary conditions, and an exact consistent tangent. Such a
formulation must recover the analytical small-strain solution in the linear
limit while providing numerical robustness for analyzing mesh convergence,
internal length scales, localized deformation regularization, material
parameter sensitivity, and higher-order boundary constraints under large
deformations.

The formulation developed in this work adopts the Green--Lagrange strain
tensor $\mathbf{E}$ and its material gradient
$\bm{\mathcal{G}}=\nabla_0\mathbf{E}$ as primary kinematic variables. The
strain-energy function consists of a classical Saint--Venant--Kirchhoff
contribution and a higher-order contribution governed by the isotropic
sixth-order Mindlin modulus tensor. A bicubic $C^1$ Bogner--Fox--Schmit (BFS)
element interpolation is employed for the primal displacement field.
Exact consistent linearization yields a tangent stiffness matrix decomposed into
classical material, classical geometric, gradient material, and gradient
geometric contributions.

The implementation is verified against the analytical benchmark of
Shekarchizadeh et al. \cite{ShekarchizadehAbaliBersani2022} in the
infinitesimal limit. The framework is subsequently applied to investigate
finite-strain mesh convergence, internal length-scale effects, localized
field regularization, sensitivity to the five constitutive gradient
constants, higher-order boundary constraints, and the progressive
partitioning of strain energy between classical and gradient mechanisms
under finite deformation.

\section{Finite-strain kinematics}
\label{sec:kinematics}

\subsection{Reference and current configurations}
\label{sec:reference-current-configurations}

Let $\mathcal{B}_0$ denote the reference configuration of the body and
$\mathcal{B}_t$ its current configuration at time $t$. A material point is
labeled by $\mathbf{X}\in\mathcal{B}_0$ and occupies the spatial position
$\mathbf{x}\in\mathcal{B}_t$. The deformation is assumed to be a smooth,
orientation-preserving mapping ensuring $J>0$ everywhere in $\mathcal{B}_0$.
Material indices are denoted by upper-case letters $(I,J,K,\ldots)$ and
spatial indices by lower-case letters $(i,j,k,\ldots)$. The material
gradient is denoted by $\nabla_0$. The reference description is retained
throughout; accordingly, the kinematic variables, internal work, and
finite-element integrations are expressed with respect to
$\mathcal{B}_0$ \cite{Holzapfel2000,Ogden1997,BonetWood2008}.

\subsection{Motion and deformation map}
\label{sec:motion-deformation-map}

The motion is written in terms of the displacement field
$\mathbf{u}(\mathbf{X},t)$ as
\begin{equation}
 \mathbf{x}=\boldsymbol{\chi}(\mathbf{X},t)
 =\mathbf{X}+\mathbf{u}(\mathbf{X},t).
 \label{eq:motion-map}
\end{equation}
The displacement field is the sole primary unknown. No independent
microrotation field is introduced.

\subsection{Deformation gradient}
\label{sec:deformation-gradient}

The deformation gradient follows from the motion as
\begin{equation}
 \mathbf{F}=\nabla_0\boldsymbol{\chi}
 =\mathbf{I}+\nabla_0\mathbf{u},
 \qquad
 F_{iI}=\frac{\partial x_i}{\partial X_I},
 \qquad
 J=\det\mathbf{F}>0.
 \label{eq:deformation-gradient}
\end{equation}
The deformation gradient $\mathbf{F}$ maps infinitesimal material line
elements to their spatial counterparts and provides the kinematic basis for
the total-Lagrangian formulation \cite{Holzapfel2000,BonetWood2008}.

\subsection{Right Cauchy--Green tensor}
\label{sec:right-cauchy-green}

The right Cauchy--Green tensor is defined by
\begin{equation}
 \mathbf{C}=\mathbf{F}^{T}\mathbf{F},
 \qquad
 C_{IJ}=F_{iI}F_{iJ}.
 \label{eq:right-cauchy-green}
\end{equation}
The tensor $\mathbf{C}$ is symmetric and invariant under superposed rigid-body
motions \cite{Holzapfel2000,Ogden1997}.

\subsection{Green--Lagrange strain}
\label{sec:green-lagrange-strain}

The Green--Lagrange strain is adopted as the classical strain variable,
\begin{equation}
 \mathbf{E}=\tfrac{1}{2}\left(\mathbf{F}^{T}\mathbf{F}-\mathbf{I}\right)
 =\tfrac{1}{2}\left(\mathbf{C}-\mathbf{I}\right).
 \label{eq:green-lagrange-strain}
\end{equation}
In components,
\begin{equation}
 E_{IJ}=\tfrac{1}{2}\left(F_{iI}F_{iJ}-\delta_{IJ}\right),
 \qquad E_{IJ}=E_{JI}.
 \label{eq:green-lagrange-strain-components}
\end{equation}
In terms of the displacement gradient \eqref{eq:deformation-gradient}, the
Green--Lagrange strain is expressed as
\begin{equation}
 \mathbf{E}=\tfrac{1}{2}\left(\nabla_0\mathbf{u}
 +(\nabla_0\mathbf{u})^{T}
 +(\nabla_0\mathbf{u})^{T}\nabla_0\mathbf{u}\right).
 \label{eq:green-lagrange-strain-displacement}
\end{equation}
The quadratic term is retained in the nonlinear formulation
\cite{Holzapfel2000,Ogden1997,SimoHughes1998}.

\subsection{Material strain-gradient tensor}
\label{sec:material-strain-gradient}

The first strain-gradient measure is the material gradient of the
Green--Lagrange strain,
\begin{equation}
 \bm{\mathcal{G}}=\nabla_0\mathbf{E},
 \qquad
 \mathcal{G}_{IJK}=E_{IJ,K}.
 \label{eq:material-strain-gradient}
\end{equation}
Its component form is
\begin{equation}
 \mathcal{G}_{IJK}
 =\tfrac{1}{2}\left(F_{iI,K}F_{iJ}+F_{iI}F_{iJ,K}\right),
 \qquad F_{iI,K}=u_{i,IK}.
 \label{eq:material-strain-gradient-components}
\end{equation}
Consequently, $\bm{\mathcal{G}}$ is symmetric with respect to its first two
indices:
\begin{equation}
 \mathcal{G}_{IJK}=\mathcal{G}_{JIK}.
 \label{eq:material-strain-gradient-symmetry}
\end{equation}
No index symmetry applies to the gradient index $K$. Adopting
$\nabla_0\mathbf{E}$ extends Mindlin Form-II gradient kinematics to finite
deformation in a referential description
\cite{MindlinEshel1968,AbaliMullerEremeyev2015}. Because $\nabla_0\mathbf{E}$
contains second material derivatives of displacement, conforming
finite-element approximations require $C^1$ continuity
\cite{ShuKingFleck1999}.

\subsection{Small-strain recovery}
\label{sec:small-strain-recovery}

For infinitesimal displacement gradients ($\|\nabla_0\mathbf{u}\|\ll 1$), the
quadratic term in \eqref{eq:green-lagrange-strain-displacement} is negligible,
and the finite-strain measures recover the linear kinematic variables:
\begin{equation}
 \mathbf{E}\longrightarrow\boldsymbol{\varepsilon}
 =\tfrac{1}{2}\left(\nabla\mathbf{u}+(\nabla\mathbf{u})^{T}\right),
 \qquad
 \nabla_0\mathbf{E}\longrightarrow\nabla\boldsymbol{\varepsilon}.
 \label{eq:kinematic-small-strain-recovery}
\end{equation}
In this limit, the distinction between material and spatial gradients
vanishes to first order. The recovered variables are therefore those of the
Mindlin Form-II model \cite{MindlinEshel1968}.

\subsection{Remarks on the adopted finite-strain variables}
\label{sec:remarks-finite-strain-variables}

The pair $(\mathbf{E},\bm{\mathcal{G}})$ constitutes the primary kinematic
variables throughout the formulation. The material gradient of the
Green--Lagrange strain is adopted as the higher-order kinematic variable,
as distinct from a spatial strain gradient or a rotation gradient. From the
polar decomposition $\mathbf{F}=\mathbf{R}\mathbf{U}$, the rotation tensor
$\mathbf{R}$ and its material gradient can be computed as kinematic
quantities, but neither enters as an independent field or a constitutive
argument \cite{Holzapfel2000,Ogden1997}.

The implementation is two-dimensional, with $I,J,K\in\{1,2\}$, and assumes
plane strain by default, where $E_{33}=0$ is enforced kinematically. A
plane-stress formulation, if required, operates as a constitutive reduction
(enforcing $S_{33}=0$, with $E_{33}$ determined constitutively) rather than
imposing $E_{33}=0$ as a kinematic constraint. These conventions are
retained throughout the constitutive response, the higher-order weak form,
and the $C^1$ displacement discretization.

\section{Constitutive formulation}
\label{sec:constitutive}

\subsection{Constitutive assumptions and state variables}
\label{sec:constitutive-assumptions}

An isothermal, elastic material response is considered. The Helmholtz
free-energy density is defined per unit reference volume
\cite{ColemanNoll1963,Holzapfel2000} and depends on the finite-strain
variables introduced in Section~\ref{sec:kinematics}:
\begin{equation}
 \mathcal{S}=\left\{\mathbf{E},\bm{\mathcal{G}}\right\},
 \qquad
 \bm{\mathcal{G}}=\nabla_0\mathbf{E}.
 \label{eq:constitutive-state}
\end{equation}
Although $\bm{\mathcal{G}}$ is obtained from the material gradient of
$\mathbf{E}$, the pair in \eqref{eq:constitutive-state} constitutes the set of
local constitutive state variables, with kinematic compatibility ensured
through the primary displacement field $\mathbf{u}$. No independent
microrotation, microdeformation, or dissipative internal variable is
introduced. This material specialization follows the strain--strain-gradient
structure of Mindlin Form-II elasticity in a total-Lagrangian setting
\cite{Mindlin1964,MindlinEshel1968,AbaliMullerEremeyev2015}.

The reference state is homogeneous and stress-free,
\begin{equation}
 \Psi(\mathbf{0},\mathbf{0})=0,
 \qquad
 \mathbf{S}(\mathbf{0},\mathbf{0})=\mathbf{0},
 \qquad
 \boldsymbol{\Sigma}(\mathbf{0},\mathbf{0})=\mathbf{0}.
 \label{eq:constitutive-reference-state}
\end{equation}
The material is assumed to be isotropic and centrosymmetric. Material
centrosymmetry excludes coupling terms bilinear in $\mathbf{E}$ and
$\bm{\mathcal{G}}$, while material isotropy restricts the gradient
strain-energy density to five independent scalar constants
\cite{Mindlin1964,MindlinEshel1968}.

\subsection{Helmholtz free-energy density}
\label{sec:helmholtz-energy}

The selected free-energy density is additively decomposed into classical
and strain-gradient contributions,
\begin{equation}
 \Psi(\mathbf{E},\bm{\mathcal{G}})
 =\Psi_{\mathrm{SVK}}(\mathbf{E})
 +\Psi_{\mathrm{g}}(\bm{\mathcal{G}}).
 \label{eq:energy-decomposition}
\end{equation}
The classical contribution is the Saint--Venant--Kirchhoff energy,
\begin{equation}
 \Psi_{\mathrm{SVK}}(\mathbf{E})
 =\frac{\lambda}{2}\left(\tr\mathbf{E}\right)^2
 +\mu\,\mathbf{E}:\mathbf{E},
 \label{eq:svk-energy}
\end{equation}
where $\lambda$ and $\mu$ are the Lam\'e constants. Because
$\mathbf{E}$ is an objective material strain measure, the energy in
\eqref{eq:svk-energy} is frame-indifferent. It is the standard
Green-elastic Saint--Venant--Kirchhoff form used in total-Lagrangian
finite-deformation formulations \cite{Holzapfel2000,Ogden1997,BonetWood2008}.

The Mindlin Form-II contribution is written in terms of the sixth-order
material modulus tensor as
\begin{equation}
 \Psi_{\mathrm{g}}(\bm{\mathcal{G}})
 =\frac{1}{2}\mathbb{D}_{ABCPQR}
 \mathcal{G}_{ABC}\mathcal{G}_{PQR}.
 \label{eq:gradient-energy}
\end{equation}
Consequently, the complete free-energy density is
\begin{equation}
 \Psi(\mathbf{E},\bm{\mathcal{G}})
 =
 \frac{\lambda}{2}\left(\tr\mathbf{E}\right)^2
 +\mu\,\mathbf{E}:\mathbf{E}
 +\frac{1}{2}\mathbb{D}_{ABCPQR}
 \mathcal{G}_{ABC}\mathcal{G}_{PQR}.
 \label{eq:complete-free-energy}
\end{equation}
The nonlinearity of \eqref{eq:complete-free-energy} is geometric: both
$\mathbf{E}$ and $\bm{\mathcal{G}}$ depend nonlinearly on the
displacement field, while the material dependence on these variables remains
quadratic. Adopting a material strain gradient preserves the
reference-configuration character of the total-Lagrangian formulation and
provides the direct finite-strain extension of the Mindlin Form-II energy
\cite{MindlinEshel1968,AbaliMullerEremeyev2015}.

\subsection{Isotropic sixth-order gradient modulus}
\label{sec:isotropic-gradient-modulus}

The strain-gradient tensor satisfies
$\mathcal{G}_{ABC}=\mathcal{G}_{BAC}$. The modulus tensor in
\eqref{eq:gradient-energy} therefore acts on third-order tensors that
are symmetric in their first two indices. Its required minor and major
symmetries are
\begin{equation}
 \mathbb{D}_{ABCPQR}
 =\mathbb{D}_{BACPQR}
 =\mathbb{D}_{ABCQPR}
 =\mathbb{D}_{PQRABC}.
 \label{eq:gradient-modulus-symmetries}
\end{equation}
For an isotropic, centrosymmetric material, the higher-order
strain-energy density is equivalently expressed in terms of five independent
scalar constants $a_1,\ldots,a_5$:
\begin{align}
 \Psi_{\mathrm{g}}(\bm{\mathcal{G}})
 ={}& a_1\,\mathcal{G}_{ABB}\mathcal{G}_{ACC}
 + a_2\,\mathcal{G}_{AAC}\mathcal{G}_{CBB}
 + a_3\,\mathcal{G}_{AAC}\mathcal{G}_{BBC}
 \nonumber\\
 &+ a_4\,\mathcal{G}_{ABC}\mathcal{G}_{ABC}
 + a_5\,\mathcal{G}_{ABC}\mathcal{G}_{CBA}.
 \label{eq:isotropic-gradient-energy}
\end{align}
Equation~\eqref{eq:isotropic-gradient-energy} defines the isotropic invariant
parameterization corresponding to $\mathbb{D}_{ABCPQR}$. The compact tensor
representation in \eqref{eq:gradient-energy} is retained in the subsequent
variational derivations to keep the work-conjugacy relations and tensor
symmetries explicit \cite{Mindlin1964,MindlinEshel1968}.

Because $\mathbf{E}$ is dimensionless and $[\bm{\mathcal{G}}]=L^{-1}$, the
physical dimensions of the constitutive parameters are
\begin{equation}
 [\lambda]=[\mu]=\mathrm{Pa},
 \qquad
 [a_1]=\cdots=[a_5]=[\mathbb{D}_{ABCPQR}]
 =\mathrm{Pa}\,\mathrm{m}^{2}.
 \label{eq:constitutive-units}
\end{equation}
Consequently, each contribution in \eqref{eq:complete-free-energy} represents
energy per unit reference volume.

\subsection{Work-conjugate stresses}
\label{sec:work-conjugate-stresses}

The stress measures work-conjugate to $\mathbf{E}$ and $\bm{\mathcal{G}}$
are the second Piola--Kirchhoff stress $\mathbf{S}$ and the material
double stress $\boldsymbol{\Sigma}$, respectively:
\begin{equation}
 \mathbf{S}=\frac{\partial\Psi}{\partial\mathbf{E}},
 \qquad
 \boldsymbol{\Sigma}=\frac{\partial\Psi}{\partial\bm{\mathcal{G}}}.
 \label{eq:work-conjugate-stresses}
\end{equation}
Differentiation of \eqref{eq:complete-free-energy} gives
\begin{equation}
 S_{IJ}=\lambda\left(\tr\mathbf{E}\right)\delta_{IJ}
 +2\mu E_{IJ}.
 \label{eq:second-piola-stress}
\end{equation}
The associated first Piola--Kirchhoff stress is
\begin{equation}
 P_{iJ}=F_{iI}S_{IJ}.
 \label{eq:first-piola-stress}
\end{equation}
The material double stress is
\begin{equation}
 \Sigma_{IJK}=\mathbb{D}_{IJKPQR}\mathcal{G}_{PQR}.
 \label{eq:double-stress}
\end{equation}
The stress symmetries follow from the symmetries of the constitutive
variables and moduli:
\begin{equation}
 S_{IJ}=S_{JI},
 \qquad
 \Sigma_{IJK}=\Sigma_{JIK}.
 \label{eq:stress-symmetries}
\end{equation}
The scalar contractions $\mathbf{S}:\delta\mathbf{E}$ and
$\boldsymbol{\Sigma}:\delta\bm{\mathcal{G}}$ represent internal virtual-work
densities per unit reference volume, providing the finite-strain
counterparts of the classical Cauchy stress and double-stress relations in
linear strain-gradient elasticity \cite{MindlinEshel1968,Germain1973}.

\subsection{Constitutive tangent tensors}
\label{sec:constitutive-tangents}

The classical and gradient constitutive tangent tensors are defined by
\begin{equation}
 \mathbb{C}_{IJKL}
 :=\frac{\partial S_{IJ}}{\partial E_{KL}}
 =\lambda\delta_{IJ}\delta_{KL}
 +\mu\left(\delta_{IK}\delta_{JL}
 +\delta_{IL}\delta_{JK}\right),
 \label{eq:classical-tangent}
\end{equation}
and
\begin{equation}
 \frac{\partial\Sigma_{IJK}}{\partial\mathcal{G}_{PQR}}
 =\mathbb{D}_{IJKPQR}.
 \label{eq:gradient-tangent}
\end{equation}
Owing to the additive decomposition of the free-energy density
\eqref{eq:complete-free-energy}, the mixed classical--gradient coupling
tangents vanish identically:
\begin{equation}
 \frac{\partial S_{IJ}}{\partial\mathcal{G}_{PQR}}=0,
 \qquad
 \frac{\partial\Sigma_{IJK}}{\partial E_{PQ}}=0.
 \label{eq:zero-mixed-tangents}
\end{equation}
The fourth-order modulus $\mathbb{C}$ possesses standard minor and major
elastic symmetries,
\begin{equation}
 \mathbb{C}_{IJKL}
 =\mathbb{C}_{JIKL}
 =\mathbb{C}_{IJLK}
 =\mathbb{C}_{KLIJ},
 \label{eq:classical-tangent-symmetries}
\end{equation}
while the sixth-order modulus $\mathbb{D}$ satisfies the symmetries in
\eqref{eq:gradient-modulus-symmetries}. These major symmetries ensure that the
material contributions to the consistent tangent stiffness matrix are
symmetric.

\subsection{Small-strain recovery}
\label{sec:constitutive-small-strain-recovery}

Recalling the kinematic reduction $\mathbf{E}\longrightarrow\boldsymbol{\varepsilon}$ and $\bm{\mathcal{G}}=\nabla_0\mathbf{E}\longrightarrow\nabla\boldsymbol{\varepsilon}$ from Section~\ref{sec:small-strain-recovery}, the complete free-energy density in \eqref{eq:complete-free-energy} reduces to
\begin{equation}
 \Psi\longrightarrow
 \frac{\lambda}{2}\left(\tr\boldsymbol{\varepsilon}\right)^2
 +\mu\,\boldsymbol{\varepsilon}:\boldsymbol{\varepsilon}
 +\frac{1}{2}\mathbb{D}_{ijkpqr}
 \varepsilon_{ij,k}\varepsilon_{pq,r},
 \label{eq:small-strain-free-energy}
\end{equation}
and the work-conjugate stresses become
\begin{equation}
 \mathbf{S}\longrightarrow\boldsymbol{\sigma},
 \qquad
 \mathbf{P}\longrightarrow\boldsymbol{\sigma},
 \qquad
 \sigma_{ij}=\lambda\left(\tr\boldsymbol{\varepsilon}\right)\delta_{ij}
 +2\mu\varepsilon_{ij},
 \qquad
 \boldsymbol{\Sigma}\longrightarrow\mathbf{h},
 \qquad
 h_{ijk}=\mathbb{D}_{ijkpqr}\varepsilon_{pq,r}.
 \label{eq:small-strain-constitutive-recovery}
\end{equation}
The constitutive formulation thus recovers the linear isotropic Mindlin
Form-II strain-gradient relations used in the analytical benchmark
\cite{Mindlin1964,MindlinEshel1968}.

\section{Variational formulation}
\label{sec:variational}

\subsection{Energy functionals and Hamilton's principle}
\label{sec:variational-principle}

The total-Lagrangian variational formulation is posed on the reference domain
$\mathcal{B}_0$. Let $\rho_0$ denote the reference mass density,
$\mathbf{f}_0$ the body force per unit reference volume, and
$\bar{\mathbf{T}}_0$ the prescribed nominal traction on the traction
boundary $\partial\mathcal{B}_{0t}$. The kinetic energy and internal
Helmholtz free energy are
\begin{equation}
 \mathcal{T}
 =\int_{\mathcal{B}_0}\frac{1}{2}\rho_0\,
 \dot{\mathbf{u}}\cdot\dot{\mathbf{u}}\,\mathrm{d}V_0,
 \qquad
 \Pi_{\mathrm{int}}
 =\int_{\mathcal{B}_0}
 \Psi\left(\mathbf{E},\bm{\mathcal{G}}\right)\,\mathrm{d}V_0.
 \label{eq:kinetic-and-internal-energy}
\end{equation}
The external virtual work is
\begin{equation}
 \delta\mathcal{W}_{\mathrm{ext}}
 =\int_{\mathcal{B}_0}\mathbf{f}_0\cdot\delta\mathbf{u}\,\mathrm{d}V_0
 +\int_{\partial\mathcal{B}_{0t}}
 \bar{\mathbf{T}}_0\cdot\delta\mathbf{u}\,\mathrm{d}A_0.
 \label{eq:external-virtual-work}
\end{equation}

For an action functional $\mathcal{A}$, Hamilton's principle is stated as
\begin{equation}
 \delta\mathcal{A}
 +\int_{t_1}^{t_2}\delta\mathcal{W}_{\mathrm{ext}}\,\mathrm{d}t=0,
 \qquad
 \mathcal{A}
 =\int_{t_1}^{t_2}
 \left(\mathcal{T}-\Pi_{\mathrm{int}}\right)\,\mathrm{d}t,
 \qquad
 \delta\mathbf{u}(t_1)=\delta\mathbf{u}(t_2)=\mathbf{0}.
 \label{eq:hamilton-principle}
\end{equation}
This is the reference-configuration form of Hamilton's principle for an
elastic continuum with higher-order internal work
\cite{Hamilton1834,SimoHughes1998,Germain1973}.

The formulation is first stated in this general dynamic setting. All
numerical examples considered subsequently are quasi-static; hence, the
inertia contribution is omitted in the finite-element computations.

Integrating the kinetic-energy variation by parts with respect to time yields the inertial virtual work:
\begin{equation}
 \delta\int_{t_1}^{t_2}\mathcal{T}\,\mathrm{d}t
 =-\int_{t_1}^{t_2}\delta\mathcal{W}_{\mathrm{iner}}\,\mathrm{d}t,
 \qquad
 \delta\mathcal{W}_{\mathrm{iner}}
 =\int_{\mathcal{B}_0}
 \rho_0\,\ddot{\mathbf{u}}\cdot\delta\mathbf{u}\,\mathrm{d}V_0.
 \label{eq:inertia-virtual-work}
\end{equation}
The resulting instantaneous variational statement is
\begin{equation}
 \delta\Pi_{\mathrm{int}}
 -\delta\mathcal{W}_{\mathrm{ext}}
 +\delta\mathcal{W}_{\mathrm{iner}}=0.
 \label{eq:instantaneous-variational-statement}
\end{equation}
In the quasi-static regime, the inertial term $\delta\mathcal{W}_{\mathrm{iner}}$ vanishes.

\subsection{Internal virtual work and kinematic variations}
\label{sec:internal-virtual-work}

Taking the first variation of the internal strain energy and substituting the work-conjugate constitutive relations from Section~\ref{sec:work-conjugate-stresses} yields the internal virtual work:
\begin{equation}
 \delta\Pi_{\mathrm{int}}
 =\int_{\mathcal{B}_0}
 \left[
 \mathbf{S}:\delta\mathbf{E}
 +\boldsymbol{\Sigma}:\delta\bm{\mathcal{G}}
 \right]\,\mathrm{d}V_0.
 \label{eq:internal-virtual-work}
\end{equation}
The variations of the kinematic measures with respect to the displacement field are
\begin{equation}
 \delta F_{iI}=\delta u_{i,I},
 \qquad
 \delta E_{IJ}
 =\frac{1}{2}\left(F_{iI}\delta u_{i,J}
 +F_{iJ}\delta u_{i,I}\right),
 \qquad
 \delta\mathcal{G}_{IJK}=\delta E_{IJ,K}.
 \label{eq:variation-chain}
\end{equation}
Using the index symmetry $S_{IJ}=S_{JI}$ yields the virtual-work identity
\begin{equation}
 S_{IJ}\,\delta E_{IJ}
 =P_{iJ}\,\delta u_{i,J},
 \qquad
 P_{iJ}=S_{IJ}F_{iI}.
 \label{eq:piola-virtual-work-identity}
\end{equation}
The second term in \eqref{eq:internal-virtual-work} involves the
variation of the material strain gradient $\delta\bm{\mathcal{G}}$ and,
consequently, second derivatives of the virtual displacement
$\delta\mathbf{u}$, representing the higher-order internal virtual work of
Mindlin Form-II gradient elasticity \cite{MindlinEshel1968,Germain1973}.

\subsection{Strong form and boundary terms}
\label{sec:strong-form-boundary-terms}

Applying the divergence theorem first to
$\boldsymbol{\Sigma}:\delta\bm{\mathcal{G}}$ and then to the remaining
term containing $\delta u_{i,J}$ gives the material hyperstress
\begin{equation}
 \mathcal{H}_{iJK}=\Sigma_{IJK}F_{iI}.
 \label{eq:material-hyperstress}
\end{equation}
The resulting momentum balance equation in the reference domain is
\begin{equation}
 P_{iJ,J}-(\Sigma_{IJK,K}F_{iI})_{,J}+f_{0i}
 =\rho_0\ddot{u}_i
 \qquad\text{in }\mathcal{B}_0.
 \label{eq:finite-strain-balance}
\end{equation}
The gradient contribution in \eqref{eq:finite-strain-balance} must
remain in the form $(\Sigma_{IJK,K}F_{iI})_{,J}$. In general, it is not
the divergence of $\mathcal{H}_{iJK}$ because differentiation of
\eqref{eq:material-hyperstress} also produces the term
$\Sigma_{IJK}F_{iI,K}$. At finite deformation, $\mathbf{F}$ is generally
nonuniform and this product-rule term does not vanish. Replacing the
operator by a double divergence of $\mathcal{H}_{iJK}$ would therefore
remove a finite-strain contribution and alter the governing balance.

The higher-order surface work follows from the identity
\begin{equation}
 \Sigma_{IJK}N_K\,\delta E_{IJ}
 =\Sigma_{IJK}F_{lI}N_K\,\delta u_{l,J}
 =\mathcal{H}_{lJK}N_K\,\delta u_{l,J},
 \label{eq:higher-order-boundary-identity}
\end{equation}
where $\mathbf{N}$ is the outward unit normal to
$\partial\mathcal{B}_0$. The effective single traction and the double traction are defined, respectively, by
\begin{equation}
 \left(P_{iJ}-\Sigma_{IJK,K}F_{iI}\right)N_J,
 \qquad
 \mathcal{D}_i=\mathcal{H}_{iJK}N_KN_J.
 \label{eq:natural-traction-operators}
\end{equation}
These are the finite-strain counterparts of the traction and double
traction in Mindlin-type first strain-gradient elasticity
\cite{MindlinEshel1968,Germain1973}.

Let $\partial\mathcal{B}_{0u}$ and
$\partial\mathcal{B}_{0t}$ denote complementary displacement and
traction portions of the boundary $\partial\mathcal{B}_0$. The essential boundary conditions prescribe the
displacement and its normal material derivative:
\begin{equation}
 u_i=\bar{u}_i,
 \qquad
 u_{i,N}:=u_{i,J}N_J=\bar{g}_i
 \qquad\text{on }\partial\mathcal{B}_{0u}.
 \label{eq:essential-boundary-conditions}
\end{equation}
The normal derivative is the additional essential boundary variable in a
displacement-based $C^1$ formulation. It arises from the higher-order
weak form induced by the dependence of the energy on
$\nabla_0\mathbf{E}$; it is not an independently introduced kinematic
field.

On $\partial\mathcal{B}_{0t}$, the natural boundary data are the single
traction and double traction in \eqref{eq:natural-traction-operators}. For
boundaries with unconstrained normal slope, the natural boundary condition is
the homogeneous double traction $\mathcal{D}_i=0$. The
higher-order boundary structure is consistent with the displacement-only
$C^1$ setting, for which the normal slope is the additional essential
quantity \cite{Toupin1962,Germain1973,ShuKingFleck1999}.

\subsection{Small-strain recovery}
\label{sec:variational-small-strain-recovery}

Under the kinematic and constitutive limits established in Sections~\ref{sec:small-strain-recovery} and~\ref{sec:constitutive-small-strain-recovery}, with $F_{iI}\longrightarrow\delta_{iI}$, $\rho_0\longrightarrow\rho$, and $f_{0i}\longrightarrow b_i$, the finite-strain balance \eqref{eq:finite-strain-balance} reduces to
\begin{equation}
 \sigma_{ij,j}-h_{ijk,jk}+b_i
 =\rho\ddot{u}_i.
 \label{eq:small-strain-balance}
\end{equation}
The natural operators reduce to
\begin{equation}
 \left(\sigma_{ij}-h_{ijk,k}\right)n_j,
 \qquad
 h_{ijk}n_kn_j,
 \label{eq:small-strain-natural-operators}
\end{equation}
and the essential data reduce to the displacement and normal-slope
conditions of linear Mindlin Form-II gradient elasticity. Consequently, the
variational formulation demonstrates complete consistency with the
analytical benchmark verified in Section~\ref{sec:benchmark-verification}.

\section{Finite element discretization}
\label{sec:finite-element-discretization}

\subsection{Displacement-based \texorpdfstring{$C^1$}{C1} approximation}
\label{sec:c1-bfs-approximation}

The displacement-based weak form in Section~\ref{sec:internal-virtual-work}
contains $\delta\bm{\mathcal{G}}=\nabla_0\delta\mathbf{E}$ and therefore
second material derivatives of the virtual displacement. The trial and
test fields must consequently possess square-integrable second
derivatives. A bicubic Bogner--Fox--Schmit (BFS) interpolation is adopted
on rectangular elements because it provides the required $C^1$
continuity across element boundaries \cite{BognerFoxSchmit1966,ShuKingFleck1999}.

Consider a rectangular element $\mathcal{B}_0^e$ with four nodes. For
each in-plane displacement component, the nodal variables are
\begin{equation}
 \left(u_i,\,u_{i,1},\,u_{i,2},\,u_{i,12}\right),
 \qquad i=1,2.
 \label{eq:bfs-nodal-variables}
\end{equation}
Thus, each node carries eight degrees of freedom and the element vector
contains 32 degrees of freedom. The implementation stores the local
variables in component blocks,
\begin{equation}
 \mathbf{d}^e=
 \left[\left(\mathbf{d}_1^e\right)^T,
 \left(\mathbf{d}_2^e\right)^T\right]^T,
 \qquad
 \mathbf{d}_i^e=
 \left[u_i^{(1)},\,u_{i,1}^{(1)},\,u_{i,2}^{(1)},\,u_{i,12}^{(1)},\,
 \ldots,\,
 u_i^{(4)},\,u_{i,1}^{(4)},\,u_{i,2}^{(4)},\,u_{i,12}^{(4)}\right]^T
 \in\mathbb{R}^{16},
 \qquad i=1,2.
 \label{eq:bfs-element-dof-vector}
\end{equation}

The displacement and its first derivatives are continuous across adjacent
elements, whereas the second derivatives are evaluated elementwise. This
is sufficient for the $H^2$ regularity required by the higher-order weak
form. A $C^0$ displacement interpolation would instead require a mixed
formulation with additional fields \cite{ShuKingFleck1999,AmanatidouAravas2002}.

\subsection{Element mapping and BFS interpolation}
\label{sec:bfs-element-mapping}

The parent coordinates $(\xi,\eta)\in[-1,1]^2$ are mapped affinely to
the rectangular reference element according to
\begin{equation}
 X_1=X_1^{\min}+\frac{h_1}{2}(\xi+1),
 \qquad
 X_2=X_2^{\min}+\frac{h_2}{2}(\eta+1),
 \qquad
 \left|\mathbf{J}_e\right|=\frac{h_1h_2}{4},
 \label{eq:bfs-affine-mapping}
\end{equation}
where $h_1$ and $h_2$ are the element dimensions in the material
coordinate directions. The coordinate mapping is affine for the
rectangular reference mesh. Accordingly,
$\partial/\partial X_1=(2/h_1)\partial/\partial\xi$ and
$\partial/\partial X_2=(2/h_2)\partial/\partial\eta$; the second
material derivatives use the corresponding products of these factors.

The one-dimensional cubic Hermite functions are
\begin{align}
 H_0^{(1)}(s)&=\frac{1}{4}(1-s)^2(2+s),
 & H_1^{(1)}(s)&=\frac{1}{4}(1-s)^2(1+s),
 \nonumber\\
 H_0^{(2)}(s)&=\frac{1}{4}(1+s)^2(2-s),
 & H_1^{(2)}(s)&=-\frac{1}{4}(1+s)^2(1-s),
 \qquad s\in[-1,1].
 \label{eq:hermite-functions}
\end{align}
For node $a$, let $(r_a,s_a)$ denote its endpoint indices in the
$(\xi,\eta)$ directions, with
$(r_a,s_a)=(1,1),(2,1),(2,2),(1,2)$ for $a=1,\ldots,4$. The four BFS
functions associated with this node are
\begin{align}
 N_a^{(1)}&=H_0^{(r_a)}(\xi)H_0^{(s_a)}(\eta),
 \nonumber\\
 N_a^{(2)}&=\frac{h_1}{2}H_1^{(r_a)}(\xi)H_0^{(s_a)}(\eta),
 \nonumber\\
 N_a^{(3)}&=\frac{h_2}{2}H_0^{(r_a)}(\xi)H_1^{(s_a)}(\eta),
 \nonumber\\
 N_a^{(4)}&=\frac{h_1h_2}{4}H_1^{(r_a)}(\xi)H_1^{(s_a)}(\eta).
 \label{eq:bfs-shape-functions}
\end{align}
The scaling in \eqref{eq:bfs-shape-functions} ensures interpolation of
the material derivatives appearing in \eqref{eq:bfs-nodal-variables}.

The element displacement field is interpolated as
\begin{align}
 u_i^h(\mathbf{X})
 =\sum_{a=1}^{4}\bigl[
 &N_a^{(1)}u_i^{(a)}
 +N_a^{(2)}u_{i,1}^{(a)}
 \nonumber\\
 &+N_a^{(3)}u_{i,2}^{(a)}
 +N_a^{(4)}u_{i,12}^{(a)}
 \bigr].
 \label{eq:bfs-displacement-interpolation}
\end{align}
Differentiation of \eqref{eq:bfs-displacement-interpolation} gives the
first and second material derivatives,
\begin{align}
 u_{i,J}^h
 &=\sum_{a=1}^{4}\bigl[
 N_{a,J}^{(1)}u_i^{(a)}+N_{a,J}^{(2)}u_{i,1}^{(a)}
 +N_{a,J}^{(3)}u_{i,2}^{(a)}+N_{a,J}^{(4)}u_{i,12}^{(a)}
 \bigr],
 \label{eq:bfs-first-derivatives}\\
 u_{i,JK}^h
 &=\sum_{a=1}^{4}\bigl[
 N_{a,JK}^{(1)}u_i^{(a)}+N_{a,JK}^{(2)}u_{i,1}^{(a)}
 +N_{a,JK}^{(3)}u_{i,2}^{(a)}+N_{a,JK}^{(4)}u_{i,12}^{(a)}
 \bigr].
 \label{eq:bfs-second-derivatives}
\end{align}

\subsection{Element kinematics and virtual strain operators}
\label{sec:element-kinematics}

For compact element expressions, let $\mathbf{N}_i^{*}$, $\mathbf{B}_{iJ}^{*}$,
and $\mathbf{G}_{iJK}^{*}$ be the 32-component assembled shape-function
vectors satisfying
\begin{equation}
 u_i^h=\left(\mathbf{N}_i^{*}\right)^T\mathbf{d}^e,
 \qquad
 u_{i,J}^h=\left(\mathbf{B}_{iJ}^{*}\right)^T\mathbf{d}^e,
 \qquad
 u_{i,JK}^h=\left(\mathbf{G}_{iJK}^{*}\right)^T\mathbf{d}^e.
 \label{eq:element-shape-function-vectors}
\end{equation}

The finite-strain kinematic variables are evaluated at each integration
point as
\begin{equation}
 F_{iJ}^h=\delta_{iJ}+u_{i,J}^h,
 \qquad
 E_{IJ}^h=\frac{1}{2}\left(F_{iI}^hF_{iJ}^h-\delta_{IJ}\right),
 \qquad
 \mathcal{G}_{IJK}^h
 =\frac{1}{2}\left(F_{iI,K}^hF_{iJ}^h
 +F_{iI}^hF_{iJ,K}^h\right),
 \label{eq:element-kinematic-variables}
\end{equation}
where $F_{iI,K}^h=u_{i,IK}^h$. The constitutive response is then
evaluated from the relations in Section~\ref{sec:work-conjugate-stresses}.

The virtual strain and virtual strain-gradient fields are expressed
in terms of the virtual nodal displacement vector $\delta\mathbf{d}^e$ as
\begin{equation}
 \delta E_{IJ}^h
 =\left(\mathbf{B}_{\delta E,IJ}^{*}\right)^T\delta\mathbf{d}^e,
 \qquad
 \mathbf{B}_{\delta E,IJ}^{*}
 =\frac{1}{2}\left(F_{iI}^h\mathbf{B}_{iJ}^{*}
 +F_{iJ}^h\mathbf{B}_{iI}^{*}\right),
 \label{eq:element-virtual-strain-operator}
\end{equation}
and
\begin{equation}
 \delta\mathcal{G}_{IJK}^h
 =\left(\mathbf{B}_{\delta\mathcal{G},IJK}^{*}\right)^T
 \delta\mathbf{d}^e,
 \label{eq:element-virtual-strain-gradient}
\end{equation}
with
\begin{align}
 \mathbf{B}_{\delta\mathcal{G},IJK}^{*}
 =\frac{1}{2}\bigl[
 &F_{iI,K}^h\mathbf{B}_{iJ}^{*}
 +F_{iI}^h\mathbf{G}_{iJK}^{*}
 \nonumber\\
 &+F_{iJ,K}^h\mathbf{B}_{iI}^{*}
 +F_{iJ}^h\mathbf{G}_{iIK}^{*}
 \bigr].
 \label{eq:element-virtual-strain-gradient-operator}
\end{align}
Equations~\eqref{eq:element-virtual-strain-operator} and
\eqref{eq:element-virtual-strain-gradient-operator} are the discrete
counterparts of the exact variation chain in
\eqref{eq:variation-chain}.

\subsection{Element residual and numerical integration}
\label{sec:element-residual-integration}

Substitution of the BFS fields into the internal virtual work gives the
element internal-force vector
\begin{align}
 \mathbf{r}_{\mathrm{int}}^e
 &=\int_{\mathcal{B}_0^e}
 \left[
 \mathbf{B}_{\delta E,IJ}^{*}S_{IJ}^h
 +\mathbf{B}_{\delta\mathcal{G},IJK}^{*}\Sigma_{IJK}^h
 \right]\,\mathrm{d}V_0
 \nonumber\\
 &=\int_{\mathcal{B}_0^e}
 \left[
 \mathbf{B}_{iJ}^{*}P_{iJ}^h
 +\mathbf{B}_{\delta\mathcal{G},IJK}^{*}\Sigma_{IJK}^h
 \right]\,\mathrm{d}V_0.
 \label{eq:element-internal-force}
\end{align}
The equivalence of the two forms follows from
\eqref{eq:piola-virtual-work-identity}. For dead body forces and nominal
tractions, the element external-force vector is
\begin{equation}
 \mathbf{f}_{\mathrm{ext}}^e
 =\int_{\mathcal{B}_0^e}\mathbf{N}_i^{*} f_{0i}\,\mathrm{d}V_0
 +\int_{\partial\mathcal{B}_{0t}^e}
 \mathbf{N}_i^{*}\bar{T}_{0i}\,\mathrm{d}A_0.
 \label{eq:element-external-force}
\end{equation}
For traction boundaries with unconstrained normal slope, the homogeneous
condition $\mathcal{D}_i=0$ is applied; hence, no separate double-traction
external-force vector is assembled. The quasi-static element residual is
consequently
\begin{equation}
 \mathbf{r}^e(\mathbf{d}^e)
 =\mathbf{r}_{\mathrm{int}}^e(\mathbf{d}^e)
 -\mathbf{f}_{\mathrm{ext}}^e.
 \label{eq:element-residual}
\end{equation}

All volume integrals are evaluated using $3\times3$ Gauss--Legendre
quadrature in the parent element,
\begin{equation}
 \int_{\mathcal{B}_0^e}q(\mathbf{X})\,\mathrm{d}V_0
 \approx\left|\mathbf{J}_e\right|
 \sum_{p=1}^{3}\sum_{q=1}^{3}
w_pw_q\,q(\xi_p,\eta_q),
 \label{eq:volume-quadrature}
\end{equation}
where $\xi_p,\eta_q\in\{-\sqrt{3/5},0,\sqrt{3/5}\}$ and
$w_p,w_q\in\{5/9,8/9,5/9\}$. The plane-strain implementation uses unit
reference thickness unless a reference thickness $t$ is supplied; in the
latter case the integration factor is $t\left|\mathbf{J}_e\right|$.
Boundary integrals in \eqref{eq:element-external-force} use three-point
one-dimensional Gauss--Legendre quadrature. These rules are applied consistently to the
internal force, external force, and subsequent element matrices
\cite{Hughes1987,Bathe2006}.

\subsection{Global assembly and essential boundary conditions}
\label{sec:global-assembly-essential-bc}

Let $\mathcal{I}^e$ map the 32 local degrees of freedom of element $e$
to the global degree-of-freedom vector $\mathbf{D}$. The global
quasi-static residual is assembled as
\begin{equation}
 \mathbf{R}(\mathbf{D})
 =\mathbb{A}_{e=1}^{N_{\mathrm{elem}}}
 \mathbf{r}^e(\mathbf{d}^e)
 =\mathbf{R}_{\mathrm{int}}(\mathbf{D})
 -\mathbf{F}_{\mathrm{ext}}
 =\mathbf{0}.
 \label{eq:global-residual}
\end{equation}
Here $\mathbb{A}$ denotes the standard finite-element assembly operator.
The resulting matrices and vectors are sparse, with each element
contributing a $32\times1$ residual vector and, after linearization, a
$32\times32$ element matrix \cite{Hughes1987,BelytschkoLiuMoran2000}.

The essential data in \eqref{eq:essential-boundary-conditions} are
imposed directly on the BFS displacement and normal-slope degrees of
freedom. After assembly, the prescribed degrees of freedom are separated
from the free degrees of freedom; the reduced system is formed on the
free set, while the reactions are recovered from the residual at the
prescribed set. This treatment enforces the displacement and slope data
without introducing an independent higher-order field.

\section{Consistent tangent matrix and nonlinear solution procedure}
\label{sec:consistent-tangent-newton}

\subsection{Consistent linearization of the element residual}
\label{sec:consistent-linearization}

The quasi-static element residual in \eqref{eq:element-residual} is
nonlinear because the Green--Lagrange strain, its material gradient, and
the associated virtual strain operators depend on the current element
vector $\mathbf{d}^e$. For an admissible increment
$\Delta\mathbf{d}^e$, its consistent linearization is
\begin{equation}
 \mathbf{r}^e\left(\mathbf{d}^e+\Delta\mathbf{d}^e\right)
 =\mathbf{r}^e\left(\mathbf{d}^e\right)
 +\mathbf{K}_T^e\Delta\mathbf{d}^e
 +o\left(\left\|\Delta\mathbf{d}^e\right\|\right),
 \label{eq:element-residual-linearization}
\end{equation}
where $\mathbf{K}_T^e$ is the consistent element tangent matrix. The
linearization follows directly from the second variation of the stored
energy and is therefore symmetric for the hyperelastic response adopted
in Section~\ref{sec:constitutive} \cite{SimoHughes1998,Holzapfel2000}.

For two independent element variations, the second variations of the
kinematic variables are
\begin{align}
 \Delta^2E_{IJ}
 &=\frac{1}{2}\left(
 \delta u_{i,I}\,\Delta u_{i,J}
 +\delta u_{i,J}\,\Delta u_{i,I}
 \right),
 \label{eq:second-variation-strain}\\
 \Delta^2\mathcal{G}_{IJK}
 &=\frac{1}{2}\bigl(
 \Delta u_{i,IK}\,\delta u_{i,J}
 +\Delta u_{i,I}\,\delta u_{i,JK}
 \nonumber\\
 &\hspace{23mm}
 +\Delta u_{i,JK}\,\delta u_{i,I}
 +\Delta u_{i,J}\,\delta u_{i,IK}
 \bigr).
 \label{eq:second-variation-strain-gradient}
\end{align}
Together with the constitutive tangents in
Section~\ref{sec:constitutive-tangents}, these expressions generate the
classical and strain-gradient material and geometric contributions to
$\mathbf{K}_T^e$.

\subsection{Element tangent matrix}
\label{sec:element-tangent-matrix}

The consistent tangent is decomposed as
\begin{equation}
 \mathbf{K}_T^e
 =\mathbf{K}_{\mathrm{mat,cl}}^e
 +\mathbf{K}_{\mathrm{geo,cl}}^e
 +\mathbf{K}_{\mathrm{mat,gr}}^e
 +\mathbf{K}_{\mathrm{geo,gr}}^e
 \in\mathbb{R}^{32\times32}.
 \label{eq:element-tangent-decomposition}
\end{equation}
The classical material contribution is
\begin{equation}
 \mathbf{K}_{\mathrm{mat,cl}}^e
 =\int_{\mathcal{B}_0^e}
 \mathbf{B}_{\delta E,IJ}^{*}\,
 \mathbb{C}_{IJKL}\,
 \left(\mathbf{B}_{\delta E,KL}^{*}\right)^T
 \mathrm{d}V_0,
 \label{eq:classical-material-tangent}
\end{equation}
where $\mathbb{C}_{IJKL}$ is given by
\eqref{eq:classical-tangent}. The classical geometric contribution,
which arises from \eqref{eq:second-variation-strain}, is
\begin{equation}
 \mathbf{K}_{\mathrm{geo,cl}}^e
 =\int_{\mathcal{B}_0^e}
 S_{IJ}\,\mathbf{B}_{iI}^{*}
 \left(\mathbf{B}_{iJ}^{*}\right)^T
 \mathrm{d}V_0.
 \label{eq:classical-geometric-tangent}
\end{equation}

The strain-gradient material contribution is
\begin{equation}
 \mathbf{K}_{\mathrm{mat,gr}}^e
 =\int_{\mathcal{B}_0^e}
 \mathbf{B}_{\delta\mathcal{G},IJK}^{*}\,
 \mathbb{D}_{IJKPQR}\,
 \left(\mathbf{B}_{\delta\mathcal{G},PQR}^{*}\right)^T
 \mathrm{d}V_0,
 \label{eq:gradient-material-tangent}
\end{equation}
whereas the gradient geometric contribution is
\begin{equation}
 \mathbf{K}_{\mathrm{geo,gr}}^e
 =\int_{\mathcal{B}_0^e}
 \Sigma_{IJK}
 \left[
 \mathbf{B}_{iJ}^{*}\left(\mathbf{G}_{iIK}^{*}\right)^T
 +\mathbf{G}_{iJK}^{*}\left(\mathbf{B}_{iI}^{*}\right)^T
 \right]\mathrm{d}V_0.
 \label{eq:gradient-geometric-tangent}
\end{equation}
The four terms in \eqref{eq:element-tangent-decomposition} are evaluated
at the same $3\times3$ Gauss points used for the element residual. The
major symmetries of $\mathbb{C}_{IJKL}$ and
$\mathbb{D}_{IJKPQR}$, together with the index symmetries of $\mathbf{S}$ and
$\boldsymbol{\Sigma}$, ensure that the element tangent stiffness matrix is
symmetric. Both classical and gradient geometric terms vanish at the
stress-free reference state.

\subsection{Global residual and tangent assembly}
\label{sec:global-tangent-assembly}

For the dead-load vector $\mathbf{F}_{\mathrm{ext}}$ and load factor
$\lambda_{\mathrm{load}}$, the global residual is evaluated as
\begin{equation}
 \mathbf{R}(\mathbf{D},\lambda_{\mathrm{load}})
 =\mathbf{R}_{\mathrm{int}}(\mathbf{D})
 -\lambda_{\mathrm{load}}\mathbf{F}_{\mathrm{ext}}.
 \label{eq:global-residual-load-factor}
\end{equation}
The global consistent tangent is assembled from the element matrices,
\begin{equation}
 \mathbf{K}_T(\mathbf{D})
 =\mathbb{A}_{e=1}^{N_{\mathrm{elem}}}
 \mathbf{K}_T^e(\mathbf{d}^e).
 \label{eq:global-tangent-assembly}
\end{equation}
The assembly follows the component-block BFS degree-of-freedom ordering
specified in \eqref{eq:bfs-element-dof-vector}. Sparse assembly is used
for both the residual and the tangent matrix.

Let the free and essential degree-of-freedom index sets be denoted by
$\mathcal{I}_f$ and $\mathcal{I}_e$, respectively. At a specified load
factor, the prescribed essential increment is
\begin{equation}
 \Delta\mathbf{D}_e
 =\lambda_{\mathrm{load}}\,\bar{\mathbf{D}}_e-\mathbf{D}_e,
 \label{eq:essential-increment}
\end{equation}
where $\bar{\mathbf{D}}_e$ contains the prescribed displacement and
normal-slope data. The reduced residual and linearized free system are
\begin{equation}
 \mathbf{R}_r
 =\mathbf{R}_f+\mathbf{K}_{fe}\Delta\mathbf{D}_e,
 \qquad
 \mathbf{K}_{ff}\Delta\mathbf{D}_f=-\mathbf{R}_r.
 \label{eq:reduced-newton-system}
\end{equation}
This elimination enforces essential data exactly; the residual at the
essential degrees of freedom is retained for reaction recovery.

\subsection{Newton--Raphson iteration and load stepping}
\label{sec:newton-load-stepping}

At Newton iteration $k$ of a load step, the tangent and residual are
evaluated at $\mathbf{D}^{(k)}$. The reduced system
\eqref{eq:reduced-newton-system} is solved for the free increment, and
the displacement vector is updated as
\begin{equation}
 \mathbf{D}^{(k+1)}
 =\mathbf{D}^{(k)}+\Delta\mathbf{D}^{(k)},
 \label{eq:newton-update}
\end{equation}
with the essential entries overwritten by their prescribed values at the
current load factor. This is the standard consistent Newton--Raphson
procedure for nonlinear finite-element equilibrium
\cite{SimoHughes1998,Bathe2006,BelytschkoLiuMoran2000}.

Convergence is evaluated on the reduced free-degree-of-freedom system.
The residual, increment, and energy measures are
\begin{equation}
 \eta_R^{(k)}
 =\frac{\|\mathbf{R}_r^{(k)}\|_2}{\|\mathbf{R}_r^{(0)}\|_2},
 \qquad
 \eta_D^{(k)}
 =\frac{\|\Delta\mathbf{D}_f^{(k)}\|_2}
 {\|\mathbf{D}_f^{(k+1)}\|_2},
 \qquad
 \eta_E^{(k)}
 =\frac{\left|\left(\Delta\mathbf{D}_f^{(k)}\right)^T
 \mathbf{R}_r^{(k)}\right|}
 {\left|\left(\Delta\mathbf{D}_f^{(0)}\right)^T
 \mathbf{R}_r^{(0)}\right|}.
 \label{eq:newton-convergence-measures}
\end{equation}
A load step is accepted when all three measures satisfy
\begin{equation}
 \eta_R^{(k)}<\mathrm{tol}_R,
 \qquad
 \eta_D^{(k)}<\mathrm{tol}_D,
 \qquad
 \eta_E^{(k)}<\mathrm{tol}_E,
 \label{eq:newton-convergence-test}
\end{equation}
with
\begin{equation}
 \mathrm{tol}_R=\mathrm{tol}_D=\mathrm{tol}_E=10^{-8},
 \qquad
 k_{\max}=20.
 \label{eq:newton-controls}
\end{equation}
The implementation also uses absolute safeguards of $10^{-12}$ for the
residual, increment, and energy measures when a normalizing denominator
vanishes or falls below numerical tolerance.

The external load is advanced incrementally. When a load step fails to
converge or is declared divergent, its increment is reduced by
\begin{equation}
 \Delta\lambda_{\mathrm{load}}\leftarrow\alpha\Delta\lambda_{\mathrm{load}},
 \qquad \alpha=0.5.
 \label{eq:load-cutback}
\end{equation}
A target step is abandoned when the permitted number of cutbacks is
exhausted or when the trial increment falls below the prescribed minimum.
When adaptive stepping is enabled, a successful step requiring no more
than $k_{\mathrm{target}}=6$ iterations permits the next increment to be
enlarged according to
\begin{equation}
 \Delta\lambda_{\mathrm{load}}\leftarrow\beta\Delta\lambda_{\mathrm{load}},
 \qquad \beta=1.2.
 \label{eq:load-growth}
\end{equation}
When a converged step requires at least $k_{\max}-1$ iterations, the
subsequent increment is reduced using the factor $\alpha$. The solution
sequence consists of residual evaluation, tangent assembly, strong
imposition of the essential data, solution of the reduced linear system,
and the convergence checks in \eqref{eq:newton-convergence-test}. The
procedure is repeated until the final load factor is reached or the
admissible load-stepping limits are encountered.

\section{Verification against the strain-gradient elasticity benchmark}

\label{sec:benchmark-verification}

\subsection{Benchmark setting}

\label{sec:benchmark-setting}

The implementation was verified against the analytical two-dimensional
strain-gradient elasticity benchmark of Shekarchizadeh et al.
\cite{ShekarchizadehAbaliBersani2022}. The problem is a simple-shear
plate of height $H=0.5\,\mathrm{mm}$ and length $L=3H$. Periodic
conditions are imposed on the lateral boundaries, with $u_y=0$ and
$u_x=u_x(y)$. The material parameters are $E=400\,\mathrm{MPa}$,
$\nu=0.49$, and $\ell\in\{0.1,0.2,0.3\}\,\mathrm{mm}$.

Two loading cases were examined. In Case~1, the lower and upper
boundaries satisfy

\begin{equation}
 u_x(0)=0,
 \qquad
 u_x(H)=\widehat{u},
 \qquad
 m_{xy}(0)=0,
 \qquad
 u_{x,y}(H)=0.
 \label{eq:benchmark-case1-boundary-conditions}
\end{equation}

In Case~2, the prescribed data are

\begin{equation}
 u_x(0)=0,
 \qquad
 t_x(H)=\widehat{t},
 \qquad
 u_{x,y}(0)=0,
 \qquad
 m_{xy}(H)=0.
 \label{eq:benchmark-case2-boundary-conditions}
\end{equation}

In \eqref{eq:benchmark-case1-boundary-conditions} and \eqref{eq:benchmark-case2-boundary-conditions}, $m_{xy}$ denotes the higher-order double-traction component conjugate to the normal slope $u_{x,y}$, which corresponds to the double-traction component $\mathcal{D}_x$ in \eqref{eq:natural-traction-operators}. Periodic boundary conditions were enforced on the lateral boundaries by coupling the corresponding nodal degrees of freedom. Zero
higher-order boundary data were applied naturally, and the specified slope
was imposed through the corresponding BFS essential degrees of freedom.

The benchmark parameters were expressed in the five-invariant notation
used here through

\begin{equation}
 a_1=2c_5,
 \qquad
 a_2=2c_3,
 \qquad
 a_3=\frac{c_4}{2},
 \qquad
 a_4=c_6,
 \qquad
 a_5=2c_7,
 \label{eq:benchmark-gradient-parameter-map}
\end{equation}

where

\begin{equation}
 c_3=c_4=\frac{\ell^2\lambda}{12},
 \qquad
 c_5=c_7=\frac{\ell^2(7\mu+3\lambda)}{120},
 \qquad
 c_6=\frac{\ell^2(7\mu-4\lambda)}{120}.
 \label{eq:benchmark-gradient-coefficients}
\end{equation}

To ensure an infinitesimal deformation regime, both the benchmark loading and the
analytical solution were scaled by a factor of $10^{-5}$. This comparison confirms
that the finite-strain formulation recovers the linear kinematics
($\mathbf{E}\to\boldsymbol{\varepsilon}$ and $\nabla_0\mathbf{E}\to\nabla\boldsymbol{\varepsilon}$)
without modifying the nonlinear solver implementation.

\subsection{Error measures and mesh sequence}

\label{sec:benchmark-error-measures}

Uniform $12\times4$, $24\times8$, and $48\times16$ BFS meshes were used,
giving 520, 1800, and 6664 global degrees of freedom. The relative
displacement, $H^1$, and strain-gradient errors were calculated as

\begin{equation}
 e_{L^2}=\frac{\|u_x^h-u_x\|_{L^2}}{\|u_x\|_{L^2}},
 \qquad
 e_{H^1}=\frac{\|u_x^h-u_x\|_{H^1}}{\|u_x\|_{H^1}},
 \qquad
 e_G=\frac{\|u_{x,yy}^h-u_{x,yy}\|_{L^2}}
 {\|u_{x,yy}\|_{L^2}}.
 \label{eq:benchmark-error-measures}
\end{equation}

Periodic constraints were condensed in the benchmark analysis, and
convergence was assessed from the reduced constraint residual.

\subsection{Verification results}

\label{sec:benchmark-results}

Figure~\ref{fig:benchmark-profiles} compares the analytical displacement
profiles with the BFS solution on the finest mesh. For all three length scales and both loading cases, the numerical symbols
are visually indistinguishable from the analytical curves. The variation in
the displacement profiles with increasing $\ell$ confirms that the
strain-gradient contribution is correctly captured.

\begin{figure}[htbp]
 \centering
 \includegraphics[width=\linewidth]{Fig01_BenchmarkProfiles.eps}
 \caption{Analytical and BFS displacement profiles for the Shekarchizadeh et al.
 simple-shear benchmark on the $48\times16$ mesh. Solid curves: analytical
 solution. Open symbols: BFS solution. Left: Case~1, prescribed top
 displacement. Right: Case~2, prescribed top traction.}
 \label{fig:benchmark-profiles}
\end{figure}

Table~\ref{tab:benchmark-results} gives the errors on the finest mesh
and the rates obtained from the two finest meshes. In both cases the
observed rates are close to four in $L^2$, three in $H^1$, and two for
the strain-gradient measure. All calculations converged, and the largest
reduced periodic-equilibrium residual norm was below $6\times10^{-12}$.

\begin{table}[htbp]
 \centering
 \caption{Finest-mesh errors and observed convergence rates for the
 Shekarchizadeh et al. simple-shear benchmark. Rates are evaluated between the $24\times8$ and
 $48\times16$ meshes.}
 \label{tab:benchmark-results}
 \small
 \begin{tabular}{llrrrrrr}
 \toprule
 Case & $\ell$ [mm] & $e_{L^2}$ & $e_{H^1}$ & $e_G$ & $r_{L^2}$ & $r_{H^1}$ & $r_G$ \\
 \midrule
 1 & 0.1 & $2.262\times10^{-6}$ & $5.310\times10^{-5}$ & $3.645\times10^{-3}$ & 3.96 & 2.97 & 1.97 \\
 1 & 0.2 & $4.517\times10^{-7}$ & $1.165\times10^{-5}$ & $9.158\times10^{-4}$ & 3.99 & 2.99 & 1.99 \\
 1 & 0.3 & $1.919\times10^{-7}$ & $5.040\times10^{-6}$ & $4.073\times10^{-4}$ & 4.00 & 3.00 & 2.00 \\
 \midrule
 2 & 0.1 & $3.302\times10^{-6}$ & $5.392\times10^{-5}$ & $3.645\times10^{-3}$ & 3.96 & 2.97 & 1.97 \\
 2 & 0.2 & $4.717\times10^{-7}$ & $1.189\times10^{-5}$ & $9.158\times10^{-4}$ & 3.99 & 2.99 & 1.99 \\
 2 & 0.3 & $1.410\times10^{-7}$ & $5.155\times10^{-6}$ & $4.073\times10^{-4}$ & 4.00 & 2.99 & 2.00 \\
 \bottomrule
 \end{tabular}
\end{table}

The error histories are plotted in Figure~\ref{fig:benchmark-convergence}.
The errors drop monotonically under refinement, and combined with the
profile agreement and the reduced-residual convergence, this confirms
recovery of the small-strain Mindlin Form-II response.

\begin{figure}[htbp]
 \centering
 \includegraphics[width=\linewidth]{Fig02_BenchmarkConvergence.eps}
 \caption{Mesh-convergence results for the Shekarchizadeh et al.
 simple-shear benchmark. Upper row: Case~1. Lower row: Case~2. Dashed lines
 indicate the reference orders $O(h^4)$, $O(h^3)$, and $O(h^2)$ for the
 $L^2$, $H^1$, and strain-gradient errors, respectively.}
 \label{fig:benchmark-convergence}
\end{figure}


\section{Numerical Results and Discussion}

\label{sec:numerical-results-discussion}

The previous section verified the formulation against the analytical
benchmark. With the small-strain response confirmed, the following
examples examine the performance under finite deformation.

Unless otherwise stated, all examples use the same total-Lagrangian
formulation, displacement-based $C^1$ Bogner--Fox--Schmit finite element,
constitutive model, and Newton--Raphson solution procedure described in the
preceding sections. Only the mesh, loading, material parameters, or
boundary conditions change from one study to another.

The numerical investigation begins with a mesh-convergence analysis to
select a suitable discretization for subsequent computations.
Subsequently, the effects of the material length scale, finite
deformation, higher-order boundary conditions, and the five isotropic
strain-gradient constants are examined. Finally, the classical and
strain-gradient energy contributions are evaluated.

\subsection{Finite-strain mesh-convergence assessment}

\label{sec:finite_strain_mesh}

Following the analytical small-strain benchmark, the numerical behavior of
the fully nonlinear solver was evaluated for an in-plane cantilever subjected
to a distributed end shear.  The reference domain has $L/H=4$, unit
thickness, $E=100$, and $\nu=0.3$.  The left boundary is fully clamped,
i.e., displacement and normal slope are prescribed, while a vertical dead
traction is applied on the right boundary.  Two loading regimes were
examined: a small traction $\bar t=10^{-3}$ and a moderate
finite-deformation traction $\bar t=10^{-1}$.  For each load level, the
classical finite-strain control ($\ell/H=0$) was compared with the
finite-strain strain-gradient response ($\ell/H=0.1$).

Uniform meshes $(N_{eX},N_{eY})=(8,2)$, $(16,4)$, $(32,8)$, and
$(64,16)$ were used.  The $64\times16$ mesh, with 8840 global BFS degrees
of freedom, was taken as the numerical reference.  The relative tip
displacement and energy errors are defined by

\begin{equation}
 e_u=\frac{|u_{\mathrm{tip}}^h-u_{\mathrm{tip}}^{\mathrm{ref}}|}
 {|u_{\mathrm{tip}}^{\mathrm{ref}}|},
 \qquad
 e_{\Pi}=\frac{|\Pi_{\mathrm{int}}^h-\Pi_{\mathrm{int}}^{\mathrm{ref}}|}
 {|\Pi_{\mathrm{int}}^{\mathrm{ref}}|}.
 \label{eq:finite_mesh_errors}
\end{equation}

Figure~\ref{fig:finite_mesh_errors} shows both error measures decreasing
monotonically under mesh refinement for both the small- and finite-load
cases. Although the classical and strain-gradient materials exhibit the
same convergence trend, the strain-gradient model generally produces
slightly smaller errors on the refined meshes.  The smallest non-reference mesh error occurs on the $32\times8$ mesh. The
$64\times16$ solution is adopted as the numerical reference against which
relative convergence is assessed.

\begin{figure}[htbp]
  \centering
  \includegraphics[width=\linewidth]{Fig03_FiniteStrainMeshErrors.eps}
  \caption{Finite-strain mesh-convergence errors relative to the
  $64\times16$ numerical reference.  Top row: small-load response.
  Bottom row: moderate finite-deformation response.  The two curves in
  each panel are the classical control ($\ell/H=0$) and the
  strain-gradient material ($\ell/H=0.1$).}
  \label{fig:finite_mesh_errors}
\end{figure}

Table~\ref{tab:finite_mesh_reference} reports the converged reference
responses. All 16 analyses converged without a load-step cutback. The
small-load cases required 14--15 total Newton iterations, whereas the
finite-load cases required 22 iterations. On the reference mesh, the
strain-gradient material reduced the cantilever tip displacement by
approximately $4.3\%$ relative to the classical finite-strain control at
both load levels. This reduction reflects the characteristic stiffening
induced by the strain-gradient contribution.

\begin{table}[htbp]
\centering
\caption{Reference-mesh ($64\times16$) responses for the finite-strain
mesh-convergence study.}
\label{tab:finite_mesh_reference}
\small
\begin{tabular}{lrrrrr}
\toprule
Load & $\ell/H$ & $u_{\mathrm{tip}}/H$ &
$\Pi_{\mathrm{int}}/(EH^2)$ & Newton iterations & Cutbacks \\
\midrule
Small  & 0   & $2.4386\times10^{-3}$ & $1.2202\times10^{-8}$ & 15 & 0 \\
Small  & 0.1 & $2.3337\times10^{-3}$ & $1.1678\times10^{-8}$ & 15 & 0 \\
Finite & 0   & $2.4307\times10^{-1}$ & $1.2143\times10^{-4}$ & 22 & 0 \\
Finite & 0.1 & $2.3266\times10^{-1}$ & $1.1624\times10^{-4}$ & 22 & 0 \\
\bottomrule
\end{tabular}
\end{table}

The steady reduction in both displacement and energy errors demonstrates
that the numerical solution becomes essentially mesh independent on the
finest discretizations. The $64\times16$ mesh is therefore adopted as the
reference discretization for the remaining finite-strain simulations.


\subsection{Influence of the internal length scale}

\label{sec:length_scale_effect}

The internal length scale governs the size-dependent response of Mindlin
Form-II strain-gradient elasticity. Its influence was examined using the
cantilever configuration described in Section~\ref{sec:finite_strain_mesh}. The geometry,
classical elastic constants, boundary conditions, integration rule, and
nonlinear solution procedure were held fixed; only the dimensionless ratio
$\ell/H$ was varied.  The classical finite-strain Saint--Venant--Kirchhoff
response ($\ell/H=0$) was used as the reference state. 

The parametric sweep was performed on the converged $64\times16$ BFS mesh
identified in the mesh-convergence study. Both the small-load and moderate
finite-deformation load levels were evaluated.  The apparent stiffness and
maximum finite strain-gradient measure were defined as

\begin{equation}
 K_{\mathrm{app}}=\frac{F_{\mathrm{end}}}{|u_{\mathrm{tip}}|},
 \qquad
 \mathcal{G}_{\max}=\max_{\mathcal{B}_0}\|\nabla_0\mathbf{E}\|,
 \label{eq:length_scale_measures}
\end{equation}

where $F_{\mathrm{end}}$ is the resultant end shear force and
$u_{\mathrm{tip}}$ is the mid-height displacement at the loaded end.
Reported quantities are normalized by their classical values at
$\ell/H=0$.

Figure~\ref{fig:length_scale_effect} shows that the apparent stiffness
increases monotonically with the internal length scale. At $\ell/H=0.1$,
the normalized stiffness reaches approximately $1.05$, representing an
increase of roughly $5\%$ over the classical finite-strain reference.
Concurrently, the normalized maximum strain gradient decreases to
approximately $0.19$ at $\ell/H=0.1$. The strain-gradient contribution
therefore stiffens the cantilever and reduces the magnitude of the peak
strain-gradient response. The associated spatial regularization is
examined in Section~\ref{sec:finite_strain_fields}.

The normalized small-load and finite-load curves are visually
indistinguishable over the investigated range of $\ell/H$, indicating that
the size-dependent stiffening mechanism operates similarly across both
loading regimes. The finite-load simulations nonetheless retain the full
geometrically nonlinear kinematics and require the consistent tangent
matrix to achieve quadratic convergence.

\begin{figure}[htbp]
  \centering
  \includegraphics[width=\linewidth]{Fig04_LengthScaleEffect.eps}
  \caption{Internal-length-scale effect for the finite-strain cantilever.
  Left: normalized apparent stiffness. Right: normalized maximum finite
  strain-gradient measure. The classical response ($\ell/H=0$) is the
  reference. Small-load and finite-load curves are both shown; their
  close agreement demonstrates that the size-effect trend is preserved
  across the investigated load range.}
  \label{fig:length_scale_effect}
\end{figure}

\begin{table}[htbp]
\centering
\caption{Normalized internal-length-scale response.  Values are
representative of the converged parametric sweep; the small- and
finite-load responses are nearly identical after normalization.}
\label{tab:length_scale_effect}
\small
\begin{tabular}{rrr}
\toprule
$\ell/H$ & $K_{\mathrm{app}}/K_{\mathrm{app}}^{\ell=0}$ &
$\mathcal{G}_{\max}/\mathcal{G}_{\max}^{\ell=0}$ \\
\midrule
0.00 & 1.000 & 1.000 \\
0.05 & 1.025 & 0.38 \\
0.10 & 1.050 & 0.19 \\
\bottomrule
\end{tabular}
\end{table}

These results isolate the influence of the intrinsic material length scale
from the underlying finite-strain response. They establish the baseline for the
full-field investigation in Section~\ref{sec:finite_strain_fields}, where geometric
nonlinearity and strain-gradient regularization interact under finite shear loading.



\subsection{Finite-strain field response and strain-gradient regularization}

\label{sec:finite_strain_fields}

The physical effect of the strain-gradient contribution was examined at
the final state of the nonlinear cantilever simulation.  The cantilever
was subjected to a distributed end shear with final traction $\bar t=0.25$,
corresponding to a normalized resultant force
$F_{\mathrm{end}}/(EHt)=2.5\times10^{-3}$.  A $32\times8$ structured BFS
mesh was used, and the load was applied incrementally with the
adaptive Newton--Raphson procedure.  The classical finite-strain response
($\ell/H=0$) was compared with the gradient responses $\ell/H=0.05$ and
$\ell/H=0.1$.

All nonlinear calculations converged without a load-step cutback. The converged
response data are summarized in Table~\ref{tab:finite_field_response}.
Increasing the internal length scale reduces the final tip displacement: the
$\ell/H=0.1$ gradient solution yields a tip displacement approximately $4.7\%$ lower
than the classical finite-strain solution. This reduction corresponds to the
apparent-stiffness increase observed in the length-scale study (Section~\ref{sec:length_scale_effect}).

\begin{table}[htbp]
\centering
\caption{Final nonlinear cantilever response at
$F_{\mathrm{end}}/(EHt)=2.5\times10^{-3}$.}
\label{tab:finite_field_response}
\small
\begin{tabular}{rrrr}
\toprule
$\ell/H$ & $u_{\mathrm{tip}}/H$ & Total Newton iterations & Cutbacks \\
\midrule
0.00 & $6.0428\times10^{-1}$ & 32 & 0 \\
0.05 & $5.8978\times10^{-1}$ & 32 & 0 \\
0.10 & $5.7607\times10^{-1}$ & 32 & 0 \\
\bottomrule
\end{tabular}
\end{table}

Figure~\ref{fig:finite_strain_fields} compares the final deformed configurations
and finite-strain fields of the classical and gradient-enhanced solutions.
The top row corresponds to the classical formulation ($\ell/H=0$), and the bottom
row to the strain-gradient formulation ($\ell/H=0.1$). The first and second columns
display the longitudinal and shear Green--Lagrange strains, $E_{11}$ and $E_{12}$,
respectively, while the third column shows the dimensionless strain-gradient
intensity $H|\mathcal{G}|$.

The $E_{11}$ field displays the expected bending distribution: the lower region of
the cantilever is in tension, and the upper region in compression. The $E_{12}$
shear strain remains dominant throughout the domain due to the applied end shear.
The primary deformation mode is preserved when the strain-gradient contribution is
incorporated, confirming that the gradient model does not introduce a new deformation
mechanism. Instead, it regularizes the spatial distribution of the strain-gradient
field while leaving the overall deformation mode essentially unchanged.

This regularization is most pronounced in the $H|\mathcal{G}|$ contours. In the
classical solution, the largest strain-gradient values concentrate near the clamped
boundary, where kinematic constraints and bending deformation interact.
For $\ell/H=0.1$, the strain-gradient field in this region is noticeably smoother
and less localized at the identical contour scale. This field comparison provides a
physical explanation for the reduced tip displacement and increased apparent stiffness:
the Mindlin Form-II gradient energy penalizes sharp strain-gradient variations and
spreads localized deformation over a wider region.

\begin{figure}[htbp]
  \centering
  \includegraphics[width=\linewidth]{Fig05_FiniteStrainFields.eps}
  \caption{Final deformed cantilever configurations and finite-strain
  fields under end shear.  Top row: classical finite-strain response,
  $\ell/H=0$.  Bottom row: finite-strain strain-gradient response,
  $\ell/H=0.1$.  Columns show $E_{11}$, $E_{12}$, and the dimensionless
  strain-gradient intensity $H|\mathcal{G}|$.  Identical contour scales
  are used in corresponding top and bottom panels for direct comparison.}
  \label{fig:finite_strain_fields}
\end{figure}

These field distributions link the size-dependent stiffening observed in
Section~\ref{sec:length_scale_effect} directly to the suppression of local
strain-gradient concentrations near the clamped boundary. The nonlinear response
thereby couples finite deformation with the intrinsic-length-scale regularization
provided by the Mindlin Form-II formulation.


\subsection{Sensitivity to the Mindlin strain-gradient constants}

\label{sec:gradient_sensitivity}

The five independent isotropic Mindlin Form-II constants govern distinct modes
of strain-gradient resistance. Their relative influence was evaluated through a
one-at-a-time logarithmic sensitivity analysis around the nonlinear cantilever
baseline with $\ell/H=0.1$. The finite end-shear load ($\bar t=0.25$) and the
$32\times8$ BFS mesh were retained. For each parameter $a_i$, two perturbed simulations
were performed,

\begin{equation}
 a_i^- = 0.5a_i, \qquad a_i^+ = 2a_i,
 \label{eq:sensitivity_perturbations}
\end{equation}

with the remaining four parameters held fixed.  All perturbed nonlinear
analyses converged.

For each response quantity $Q_j$, the central logarithmic sensitivity was
computed as

\begin{equation}
 S_{ij} =
 \frac{\ln Q_j(a_i^+) - \ln Q_j(a_i^-)}
 {\ln(a_i^+/a_i^-)}.
 \label{eq:log_sensitivity}
\end{equation}

The response vector was

\begin{equation}
 \mathbf{Q}=
 \left[K_{\mathrm{app}},\;u_{\mathrm{tip}},\;
 \Pi_{\mathrm{int}},\;\max|\mathcal{G}|\right].
 \label{eq:sensitivity_responses}
\end{equation}

Positive values indicate that a response quantity increases with the corresponding
constitutive parameter, whereas negative values indicate a decrease.

Figure~\ref{fig:gradient_sensitivity} displays the resulting sensitivity matrix.
The global structural metrics---apparent stiffness, tip displacement, and stored
energy---exhibit weak sensitivity to individual parameter perturbations over the
investigated range. In contrast, the local strain-gradient intensity exhibits
pronounced sensitivity. The maximum strain gradient is most sensitive to $a_5$,
with $\partial\ln(\max|\mathcal{G}|)/\partial\ln a_5=-0.31$, followed by $a_1$ with
a sensitivity of $-0.10$. Increasing $a_5$ or $a_1$ therefore suppresses the peak
strain-gradient intensity in the deformation field. The remaining constants exert
a markedly weaker influence: $a_2$ yields a slight positive sensitivity for the
maximum gradient, whereas $a_3$ and $a_4$ exhibit weak negative sensitivities.

The sensitivity results establish a clear distinction between global structural
behavior and local material response. A parameter can have negligible influence
on the overall load--displacement response and still significantly affect the
local gradient field. This distinction is important for parameter calibration and
model reduction: the global cantilever stiffness alone cannot identify all five
Mindlin constants, whereas field-based measurements or full-field numerical
simulations provide substantially richer information for identifying $a_5$ and $a_1$.

\begin{figure}[htbp]
  \centering
  \includegraphics[width=0.82\linewidth]{Fig06_GradientSensitivity.eps}
  \caption{Logarithmic one-at-a-time sensitivity matrix for the nonlinear
  gradient cantilever at $\ell/H=0.1$.  Rows correspond to perturbed
  gradient constants and columns to response quantities.  Blue values
  denote a decrease in the response when the parameter is increased;
  red values denote an increase.  Numerical values are shown in the cells.}
  \label{fig:gradient_sensitivity}
\end{figure}

\subsection{Influence of higher-order boundary conditions}

\label{sec:higher_order_bc}

The strain-gradient formulation accommodates boundary conditions that have no
counterpart in classical elasticity. In particular, the normal displacement slope
can be prescribed as an essential boundary condition or left unconstrained with its
conjugate double traction vanishing naturally. The influence of this distinction was
evaluated for the finite-strain cantilever at $\ell/H=0.1$ on the $32\times8$ BFS mesh.

Two support conditions were compared on the clamped edge $\Gamma_{\mathrm{left}}\subset\partial\mathcal{B}_{0u}$.  The fully clamped condition imposes

\begin{equation}
 \mathbf{u}=\mathbf{0},
 \qquad
 u_{i,N}=0
 \qquad \text{on } \Gamma_{\mathrm{left}},
 \label{eq:clamped_higher_order_bc}
\end{equation}

whereas the soft-support condition imposes displacement only,

\begin{equation}
 \mathbf{u}=\mathbf{0},
 \qquad
 \mathcal{D}_i=0
 \qquad \text{on } \Gamma_{\mathrm{left}}.
 \label{eq:soft_higher_order_bc}
\end{equation}

Under the soft-support condition, the normal slope remains unconstrained and the
zero double traction is satisfied naturally. Both small-load and finite-deformation
end-shear regimes were considered.

Table~\ref{tab:higher_order_bc_response} shows that the soft-support boundary is
more compliant than the fully clamped boundary. Under finite deformation, the
soft-support tip displacement is approximately $3.6\%$ larger than the clamped
counterpart. The local effect is even more pronounced: the maximum finite
strain-gradient intensity is roughly twice the fully clamped value. A consistent
trend is observed in the small-load response.

\begin{table}[htbp]
\centering
\caption{Effect of higher-order support condition at $\ell/H=0.1$.}
\label{tab:higher_order_bc_response}
\small
\begin{tabular}{lrrrr}
\toprule
Load & Support & $u_{\mathrm{tip}}/H$ & $\max|\mathcal{G}|$ & Total iterations \\
\midrule
Small & Fully clamped & $2.3483\times10^{-3}$ & $4.3000\times10^{-3}$ & 14 \\
Small & Soft support  & $2.4345\times10^{-3}$ & $7.6645\times10^{-3}$ & 14 \\
\midrule
Finite & Fully clamped & $5.7607\times10^{-1}$ & $1.0680$ & 32 \\
Finite & Soft support  & $5.9678\times10^{-1}$ & $2.1253$ & 32 \\
\bottomrule
\end{tabular}
\end{table}

Figure~\ref{fig:higher_order_bc} illustrates the spatial distribution of the
boundary-condition effect. Panels~(a) and~(b) show the dimensionless finite
strain-gradient intensity $H|\mathcal{G}|$ in the deformed configuration for the
fully clamped and soft-support cases, respectively. Panel~(c) displays the
difference field

\begin{equation}
 \Delta(H|\mathcal{G}|)=
 H|\mathcal{G}|_{\mathrm{soft}}-
 H|\mathcal{G}|_{\mathrm{clamped}}.
 \label{eq:higher_order_bc_difference}
\end{equation}

Positive values in the difference field mark regions where the soft support
produces a larger strain gradient. The field difference is concentrated
near the left support, showing that the higher-order boundary condition
primarily influences the localized strain-gradient field while leaving the
overall deformation mode essentially unchanged.

\begin{figure}[htbp]
  \centering
  \includegraphics[width=\linewidth]{Fig07_HigherOrderBCComparison.eps}
  \caption{Higher-order boundary-condition effect at the final finite
  end-shear load.  (a) Fully clamped support:
  $\mathbf{u}=\mathbf{0}$ and
  $u_{i,N}=0$.  (b) Soft support:
  $\mathbf{u}=\mathbf{0}$ with zero natural double traction.
  (c) Difference field
  $H|\mathcal{G}|_{\mathrm{soft}}-H|\mathcal{G}|_{\mathrm{clamped}}$.
  The higher-order boundary-condition influence is concentrated near the
  left support.}
  \label{fig:higher_order_bc}
\end{figure}

Specifying displacement alone is insufficient to characterize a support condition
in strain-gradient elasticity. The choice between a prescribed normal slope and a
natural double traction influences both the local boundary-layer gradient intensity
and the global structural compliance, which the $C^1$ BFS degrees of freedom
are formulated to accommodate.

\subsection{Energetic origin of strain-gradient stiffening}

\label{sec:energy_partition}

The preceding parametric and full-field studies demonstrated that the
gradient-enhanced cantilever exhibits increased apparent stiffness and suppressed
strain-gradient localization. To trace the energetic origin of this response,
the final nonlinear solutions were decomposed according to the Helmholtz free-energy
functional,

\begin{equation}
 \Pi_{\mathrm{int}}=
 \Pi_{\mathrm{SVK}}+\Pi_{\mathrm{g}}
 =\int_{\mathcal{B}_0}\left[
 \frac{\lambda}{2}(\operatorname{tr}\mathbf{E})^2+
 \mu\,\mathbf{E}:\mathbf{E}
 \right]\,\mathrm{d}V_0
 +\int_{\mathcal{B}_0}\frac{1}{2}\,
 \mathbb{D}_{ABCPQR}\mathcal{G}_{ABC}\mathcal{G}_{PQR}\,\mathrm{d}V_0.
 \label{eq:energy_partition}
\end{equation}

The decomposition was evaluated at the final nonlinear end-shear load,
using the same $32\times8$ mesh and length-scale values as the finite-field
comparison.

Figure~\ref{fig:energy_partition} gives the normalized stored-energy
partition.  The blue component is the classical Saint--Venant--Kirchhoff
contribution, the orange component the Mindlin strain-gradient
contribution.  The dashed line is the total classical reference energy at
$\ell/H=0$.

At $\ell/H=0$, the gradient contribution vanishes. As the internal length scale
grows, the gradient-energy fraction rises from $0\%$ to approximately $1.3\%$ at
$\ell/H=0.05$ and $2.3\%$ at $\ell/H=0.1$. Although the strain-gradient contribution
represents only a small fraction of the total stored energy, it is consistent with
the enhanced stiffness and reduced strain-gradient concentration observed in the
preceding field analyses. This energetic redistribution also aligns with the lower
tip displacement and higher apparent stiffness while preserving the overall
deformation mode.

The total normalized energy decreases slightly with increasing $\ell/H$ under the
prescribed end shear. This decrease does not signify an energetic instability.
Under force-controlled loading, the stiffer gradient-enhanced cantilever undergoes
smaller displacements, so that the decrease in classical deformation energy exceeds
the positive strain-gradient energy contribution. Figure~\ref{fig:energy_partition}
separates the two coupled mechanisms: the orange component represents the direct
energetic cost of strain gradients, whereas the decrease in the blue component reflects
the suppression of macroscopic compliance and localized strain concentration.

\begin{figure}[htbp]
  \centering
  \includegraphics[width=0.78\linewidth]{Fig08_EnergyPartition.eps}
  \caption{Stored-energy partition at the final nonlinear load.  Blue bars
  are the classical Saint--Venant--Kirchhoff energy and orange bars are the
  Mindlin strain-gradient energy, both normalized by the total classical
  stored energy at $\ell/H=0$.  Labels above the bars give the fraction of
  the total energy contributed by the gradient term.}
  \label{fig:energy_partition}
\end{figure}

The energy decomposition provides thermodynamic insight into the stiffening and
regularization observed in the kinematic and field responses. The strain-gradient
contribution is not an empirical numerical penalty; it represents a physically grounded
component of the free-energy functional that scales systematically with the internal
length scale.


\section{Conclusions}
\label{sec:conclusions}

This work developed a displacement-based total-Lagrangian finite element
formulation for geometrically nonlinear Mindlin Form-II strain-gradient elasticity
under plane-strain conditions. The mechanical response was formulated using the
Green--Lagrange strain and its material gradient, combining a
Saint--Venant--Kirchhoff hyperelastic potential with an isotropic sixth-order
Mindlin gradient energy, discretized via a $C^1$ Bogner--Fox--Schmit
rectangular element.

The main findings are as follows:
\begin{itemize}
 \item The exact consistent linearization of the finite-strain residual yields
 classical material, classical geometric, gradient material, and gradient
 geometric tangent contributions, enabling quadratic convergence in a
 displacement-only Newton--Raphson solver.

 \item The analytical Shekarchizadeh--Abali--Bersani benchmark was
 recovered in the infinitesimal regime for both displacement- and
 traction-controlled cases. The observed mesh-convergence rates were
 close to four in the displacement $L^2$ measure, three in the $H^1$
 measure, and two in the strain-gradient measure.

 \item The finite-strain mesh study showed monotonic convergence of tip
 displacement and internal energy for both the classical and
gradient-enhanced materials. The refined $64\times16$ mesh served as the numerical reference for the
mesh-convergence assessment.

 \item Increasing the internal length scale enhances apparent structural
 stiffness and suppresses local strain-gradient intensity. At $\ell/H=0.1$,
 the apparent stiffness is approximately $5\%$ above the classical baseline,
 while the maximum strain-gradient measure decreases to approximately $0.19$
 of its classical reference value.

 \item Under finite end-shear loading, the tip displacement decreased from
 $0.6043H$ in the classical case to $0.5761H$ for $\ell/H=0.1$.
 Full-field contours confirm that the primary bending and shear
 distributions are preserved, while localized strain gradients are
 regularized near the clamped support.

 \item The local strain-gradient response was most sensitive to $a_5$,
 followed by $a_1$, whereas the global stiffness, tip displacement, and
 stored energy were only weakly affected by one-at-a-time changes in the
 individual gradient constants.

 \item Higher-order boundary conditions noticeably influence both the
 localized strain-gradient field and the overall structural compliance.
 Releasing the normal slope and enforcing zero natural double traction
 increased the finite-load tip displacement by approximately $3.6\%$ and
 produced a twofold increase in the peak strain-gradient intensity. In
 addition, the strain-gradient energy fraction rose systematically from zero
 in the classical case to approximately $1.3\%$ and $2.3\%$ at $\ell/H=0.05$
 and $0.1$, respectively.
\end{itemize}

The formulation is directly applicable to parametric studies of rectangular
plane-strain strips and cantilever members in which a support or loaded edge
produces localized strain-gradient concentrations. In such problems, it provides
a computational tool to evaluate the influence of the internal length scale
and the choice of higher-order boundary data. Extension beyond the moderate
deformation range considered here would require an alternative hyperelastic
strain-energy function, while three-dimensional implementations require a
corresponding $C^1$ volume discretization and calibration of the gradient constants.

\noindent\textbf{Data availability statement}

The author confirms that the data supporting the findings of this study
are available within the article.

\noindent\textbf{Conflict of interest disclosure statement}

No potential conflict of interest was reported by the author.

\section*{Acknowledgments}
The author gratefully acknowledges the financial support provided by the Haryana State Higher Education Council (HSHEC), Government of Haryana, under the Haryana State Research Fund (HSRF) scheme (Sanction Order No.: 14/66-2025 SPD/HSHEC/HSRF).


\bibliographystyle{elsarticle-num}

\bibliography{references_paper}
\end{document}